\section{Evaluation methods}\label{sec:methods}
\subsection{Experimental cases and reported quantities}
% E1: comparison inputs
Seven one-bed, no-intercooler cases from the 2017 NCCC C campaign provide the experimental comparison \cite{Morgan2020}.
The modeled bed has diameter \qty{0.64}{m}, height \qty{6}{m} and MellapakPlus 252Y packing, with specific area \qty{250}{m^{-1}} and void fraction 0.970.
\Cref{tab:nccc-one-bed-scope} gives the reported operating inputs; the missing lean-inlet temperatures in Cases 1C--3C are set to \qty{318.15}{K}.
In the nine-species model, reported dry gas mass flow and dry \COtwo{} fraction are converted to a wet feed by PC-SAFT water saturation at the gas-inlet temperature and pressure, with the remaining dry gas distributed between nitrogen and oxygen using $y_{\Otwo}/y_{\Ntwo}=0.08485753604$.
These inputs apply to the nine-species comparison and to the Case 3C reference used below; the twelve-state input basis is given in \cref{subsec:twelve-state}.

\begin{table}[pos=htbp]
  \centering
  \footnotesize
  \caption{Reported and assumed inputs and reported capture for the seven NCCC one-bed, no-intercooler cases \cite{Morgan2020}.}
  \label{tab:nccc-one-bed-scope}
  \begin{threeparttable}
  \setlength{\tabcolsep}{2.5pt}
  \begin{tabular*}{\textwidth}{@{\extracolsep{\fill}}ll
    S[table-format=5.0]
    S[table-format=4.0]
    S[table-format=2.1,table-space-text-post=\tnote{a}]
    S[table-format=2.1]
    S[table-format=3.1]
    S[table-format=1.2]
    S[table-format=1.2]
    S[table-format=1.3]
    S[table-format=2.1]@{}}
    \toprule
    Year & Case &
    \multicolumn{1}{c}{\shortstack{Lean\\flow\\(\unit{kg\,h^{-1}})}} &
    \multicolumn{1}{c}{\shortstack{Dry gas\\flow\\(\unit{kg\,h^{-1}})}} &
    \multicolumn{1}{c}{\shortstack{Lean\\temp.\\(\unit{\degreeCelsius})}} &
    \multicolumn{1}{c}{\shortstack{Gas\\temp.\\(\unit{\degreeCelsius})}} &
    \multicolumn{1}{c}{\shortstack{Pressure\\(\unit{kPa})}} &
    \multicolumn{1}{c}{\shortstack{Lean loading\\(mol CO$_2$/\\mol MEA)}} &
    \multicolumn{1}{c}{\shortstack{MEA mass\\fraction$^{\mathrm{b}}$}} &
    \multicolumn{1}{c}{\shortstack{Dry inlet\\$y_{\ce{CO2}}$}} &
    \multicolumn{1}{c}{\shortstack{Reported\\capture\\(\unit{\percent})}} \\
    \midrule
    2017 & 1C & 6185 & 1997 & 45.0\tnote{a} & 43.3 & 109.1 & 0.15 & 0.30 & 0.077 & 97.1 \\
    2017 & 2C & 7765 & 2499 & 45.0\tnote{a} & 43.1 & 109.6 & 0.20 & 0.30 & 0.077 & 92.3 \\
    2017 & 3C & 7517 & 2013 & 45.0\tnote{a} & 43.6 & 109.5 & 0.25 & 0.30 & 0.093 & 89.5 \\
    2017 & 4C & 6160 & 1500 & 44.9 & 43.3 & 107.7 & 0.25 & 0.30 & 0.108 & 88.9 \\
    2017 & 5C & 5237 & 1498 & 45.0 & 43.3 & 108.4 & 0.26 & 0.30 & 0.077 & 86.4 \\
    2017 & 6C & 7665 & 2700 & 44.8 & 43.4 & 109.7 & 0.31 & 0.31 & 0.077 & 60.2 \\
    2017 & 7C & 5414 & 1000 & 44.7 & 43.3 & 107.3 & 0.34 & 0.29 & 0.100 & 76.4 \\
    \bottomrule
  \end{tabular*}
  \begin{tablenotes}[flushleft]
    \footnotesize
    \item[] All seven calculations use reported dry gas mass and dry \(\COtwo\) fraction. Saturation at the gas-inlet temperature reconstructs the wet feed, giving \(y_{\COtwo}=0.085394\) from the reported dry value 0.093 for Case 3C.
    \item[a] Lean solvent inlet temperature was blank in the source table and was set to \qty{45.0}{\degreeCelsius} for the model input.
    \item[b] MEA mass fraction is on the CO$_2$-free solvent basis. Pressure is the reported value used throughout the bed.
  \end{tablenotes}
  \end{threeparttable}
\end{table}


% E2: observables
Capture errors and axial temperature comparisons describe different aspects of absorber performance.
For case $k$, capture is the percentage reduction in vapor \COtwo{} flow; signed error is prediction minus observation:
\begin{align}
 \eta_k&=100\left(1-\frac{F_{\COtwo,k}^{\vap}(H)}{F_{\COtwo,k}^{\vap}(0)}\right),\label{eq:capture}\\
 e_{\eta,k}&=\eta_k-\eta_k^{\mathrm{obs}},&
 \mathrm{MAE}_{\eta}&=\frac{1}{N_c}\sum_{k=1}^{N_c}|e_{\eta,k}|.\label{eq:capture-error}
\end{align}
Capture is in percent, errors are in percentage points, and the mean absolute error weights the $N_c=7$ cases equally.
Packing temperatures from Morgan et al.\ \cite{Morgan2020}, Table C2, are converted from Celsius to kelvin and from their top-to-bottom coordinate $x$ to normalized height $\zeta=1-x=z/H$.
Both predicted phase temperatures are compared with the taps because the source does not designate a measured phase.
Temperature agreement is summarized by the root-mean-square deviation over the 35 taps of the seven cases.
The sampled nine-species peak is the maximum liquid temperature at 101 equally spaced exported positions.

\subsection{Numerical solution and verification}
% E3: collocation method
Adaptive residual-controlled collocation in SciPy solves the mixed-boundary-value problem (BVP) in \cref{eq:state-derivative,eq:boundary-residual}.
The piecewise cubic approximation satisfies collocation conditions, with differential residuals assessed over each interval; it does not satisfy the differential equations exactly everywhere.
Both inlet boundaries are enforced simultaneously, and intervals with excessive residuals receive additional nodes.
The state and equation scales are obtained from the initial profile estimates, with temperature scales of \qty{400}{K}.

% E4a: derivatives
The assembled Jacobian includes the thermodynamic, transport and nonisothermal dependencies of the column equations.
CppAD supplies automatic derivatives of the thermodynamic calculation, which are combined with differentiated transport and energy expressions, including the bulk-fugacity/free-concentration film relation.
Repeated identical thermodynamic states reuse their values and derivatives.
Scaled directional comparisons with centered finite differences at lean and rich states satisfy relative and absolute tolerances of \num{2e-5} and \num{2e-8}; finite differences are used only to verify the derivatives.

% E3b: twelve-state solution
The twelve-state cell equations are solved simultaneously by IPOPT \cite{Wachter2006} through CasADi \cite{anderssonCasADiSoftwareFramework2019}, with the MA57 linear solver, an upwind conservative discretization and a scaled tolerance of \num{1e-7} within 60 iterations.
Each case is solved on 17, 33, 65 and 129 nodes clustered toward both ends, and each mesh starts from the solution on the previous one.
A fresh 17-node start whose interface-loading root search fails on the bracket $[0,1]$ is repeated once on $[0,0.5]$ with otherwise unchanged inputs.
The wall limit is \qty{3600}{s} per solve and \qty{7200}{s} at 129 nodes.
Mesh dependence is judged from the change in capture between the 65- and 129-node solutions.

% E4b: distinct initial guesses
An initial axial column profile and an initial chemical composition at each position serve different numerical purposes.
The coarse Case 3C calculation starts from a converged same-case Henry-law profile; its reactive profile initializes the refined calculation.
The equilibrium solve is initialized from conservative species amounts generated with a reaction-extent fraction of 0.001 (0.0001 for Case 1C); this affects only the initial guess.
Nearby thermodynamic states may use previously calculated conservative amounts, subject to the same equilibrium and liquid-root checks.

% E5: paired settings
Case 3C is solved at two jointly varied mesh and tolerance settings with unchanged physical equations, parameters and inlet conditions.
The calculations use 21 and 41 initial nodes with collocation tolerances 0.5 and 0.05, respectively, boundary tolerance 0.001 and a 1000-node limit.
Root-mean-square (RMS) differential residuals are normalized componentwise by $1+|f_i|$ in the scaled equations and summarized by the maximum over mesh intervals.
Outer mesh iterations, residuals at each iteration, Jacobian evaluations and equilibrium solves are reported as measures of numerical effort; mesh iterations do not count inner Newton steps.

% E6a: boundary inclusion criterion
Reported solutions must satisfy the solver stopping tests, positive species concentrations, capture between zero and 100 percent, and the postprocessed boundary-error criterion
\begin{align}
 \epsilon_{\mathrm{BC}}&=\left[\sum_{j=1}^{6}
 \left(100\frac{|q_j^{\mathrm{calc}}-q_j^{\mathrm{spec}}|}{q_j^{\mathrm{spec}}}\right)^2\right]^{1/2}\leq1.0,\label{eq:postprocessed-bc}\\
 \mathbf q&=\left(F_{\COtwo}^{\liq}(H),F_{\HtwoO}^{\liq}(H),T^{\liq}(H),
 F_{\COtwo}^{\vap}(0),F_{\HtwoO}^{\vap}(0),T^{\vap}(0)\right).
\end{align}
The norm uses six percentage-relative errors with positive specified values and temperatures in kelvin.
It differs from the solver's scaled boundary residual and RMS differential residual; pressure is enforced by the boundary equations but excluded from this postprocessed norm.
A calculation in which the reactive equilibrium cannot be solved is reported as unconverged; no substitute value is assigned.
Every nine-species calculation reported in \cref{sec:results} satisfies these criteria; grouped convergence diagnostics are listed in \cref{tab:run-diagnostics}, and \cref{tab:reactive-numerical-verification} presents the paired Case 3C refinement.
Every twelve-state solution reported there meets the IPOPT stopping test at the scaled tolerance above.

% E6b: balance diagnostics
Component and charge discrepancies quantify departures from the conservation and electroneutrality equations at each position, while axial ranges of $F_i^\vap-F_i^\liq$ test column component conservation.
The analogous energy diagnostic is
\begin{equation}
 \Delta\dot H_{\mathrm{net}}=
 \max_{\zeta}\left(\dot H^\vap-\dot H^\liq\right)
 -\min_{\zeta}\left(\dot H^\vap-\dot H^\liq\right),\label{eq:energy-range}
\end{equation}
evaluated over the 101 exported positions in watts; the additional Case 3C refinements also use 10,001 positions.
The continuous equations predict zero variation; the nonzero numerical range is reported as a diagnostic, separately from caloric-model accuracy.

\subsection{Thermodynamic, transport and operating perturbations}
% E7: thermodynamic design
One-at-a-time thermodynamic perturbations of the nine-species model test two binary interactions and two reaction constants around refined Case 3C.
The MEA--water and \COtwo--water $k_{ij}(T)$ laws, with $k_{ij}(\qty{313.15}{K})=\num{-0.0883119}$ and \num{-0.0464169}, are separately multiplied by 0.95 and 1.05, including their temperature slopes.
The scaled laws enter the liquid state and the reduced vapor \COtwo{} coefficient; the wet-gas vapor has no binary corrections to scale.
Scaling a negative coefficient above one makes it more negative.
For carbamate hydrolysis (R4) and protonated-MEA dissociation (R5), $\ln K_r^*(T)=\ln K_r(T)+\ln s$ with the same multipliers, preserving $d\ln K_r/dT$.
Each calculation uses 41 initial nodes, tolerance 0.05, boundary tolerance 0.001 and the refined reference profile as its initial guess.
Other thermodynamic, kinetic, transport and operating inputs remain fixed.

% E8: transport design
Separate multipliers of 0.9 and 1.1 test liquid viscosity, free molecular \COtwo{} diffusivity and the liquid-side \COtwo{} transfer coefficient on the same refined reference.
Viscosity changes propagate through ionic diffusivity, holdup and transfer correlations; the molecular \COtwo{} and MEA diffusivity formulas are not directly viscosity-dependent.
The diffusivity and transfer-coefficient changes also propagate through Hatta number, enhancement and combined film resistance in both the column equations and Jacobian.
Vapor properties are not directly scaled, and all remaining model inputs are unchanged.
A factor-one calculation checks equivalence with the reference, and the lower-diffusivity case uses conservative species amounts obtained by continuation from a nearby temperature.

% E9: operating design
Discrete Case 3C perturbations vary solvent circulation, lean loading and inlet temperature separately.
The liquid-to-reported-dry-gas mass ratio $L/G$ changes by $\pm10\%$ through liquid flow at fixed dry gas flow; loading changes from 0.25 to 0.225 or 0.275 mol \COtwo{} per mol MEA, and liquid-inlet temperature from \qty{318.15}{K} to \qty{313.15}{K} or \qty{323.15}{K}.
These calculations use the coarse campaign settings of 21 initial nodes, tolerance 0.5 and boundary tolerance 0.001, with a 1000-node limit.
Each condition is initialized from a Henry-law profile.
That initialization did not converge for the higher flow and the colder inlet, and both were restarted from the reactive Case 3C profile with unchanged equations and tolerances.
The lower-flow condition converged from both starts with capture differing by \num{4e-6} percentage points, and its restarted solution is reported.

\paragraph{Computational-cost measurement.}
% E10: timing protocol
Three sequential fresh-initialization Case 3C runs characterize cost at the coarse mesh and tolerance settings.
The shared Linux machine has an AMD Ryzen 9 5900XT processor, 16 physical cores and approximately 31 GB RAM; the runs use CPython 3.13.13, NumPy 2.4.4 and SciPy 1.17.1, with BLAS, OpenMP and MKL thread settings fixed to one.
BVP wall and process CPU times include the solver call and conversion of its returned solution; total wall time also includes Henry initialization, phase-profile evaluation and CSV export, but excludes program start-up, input loading and output writing.
Peak resident memory includes imported libraries.
